%! TEX root = ../main.tex

\clearpage

\section{Kiến Trúc Hệ Thống}

\subsection{Thiết Kế Kiến Trúc}
\begin{itemize}
  \item Thiết kế kiến trúc: Đưa ra một sơ đồ hiển thị các hệ thống con chính và các kho dữ liệu 
cũng như các mối liên hệ giữa các thành phần đó. Nếu thấy cần thiết, ta có thể giải thích 
thêm ở dạng văn bản: định danh từng hệ thống con chính và vai trò hay trách nhiệm của 
chúng; mô tả cách thức các hệ thống con cộng tác với nhau để đạt được chức năng mong 
đợi; không đi quá chi tiết vào từng hệ thống con. Mục đích chính là đạt được một sự hiểu 
biết chung về cách thức và lý do hệ thống được phân rã cũng như cách thức các thành 
phần riêng lẻ làm việc cùng nhau. 
  \item Những lựa chọn kiến trúc: Thảo luận về những kiến trúc khác cũng được xem xét. Ta nên 
giải thích được tại sao lại không chọn chúng.
\end{itemize}

\subsection{Mô Tả Sự Phân Rã}
\noindent <Mô tả sự phân rã của các hệ thống con trong thiết kế kiến trúc. Ta có thể giải thích thêm ở dạng 
văn bản nếu thấy cần thiết. Ta có thể lựa chọn cách mô tả theo chức năng hoặc mô tả theo hướng 
đối tượng. Đối với mô tả chức năng, sử dụng lưu đồ dòng dữ liệu mức cao và các sơ đồ phân rã 
cấu trúc. Đối với mô tả theo hướng đối tượng, sử dụng mô hình hệ thống con, các sơ đồ đối 
tượng, sơ đồ phân cấp tổng quát hóa (nếu có), sơ đồ phân cấp kết hợp (nếu có), các đặc tả giao 
diện và các sơ đồ tuần tự. >

\subsection{Cơ Sở Thiết Kế}
<Thảo luận lý do cơ bản cho việc chọn lựa kiến trúc được mô tả trong mục 3.1 bao gồm các vấn 
đề then chốt và các thỏa hiệp.> 